% !TEX program = xelatex
\documentclass[11pt]{article}
\usepackage[margin=1.2in]{geometry}
\usepackage{amsmath, amssymb, graphicx, ragged2e, titlesec, fontspec}
\setmainfont{Palatino}
\titleformat{\section}{\Large\itshape}{}{0pt}{}
\titleformat{\subsection}{\large\itshape}{}{0pt}{}
\RaggedRight \setlength{\parindent}{1em}

\title{A Note on Sawtooth Functions}
\author{A. Author}

\begin{document}
\maketitle

\section{Introduction}
We study the periodic sawtooth function and its Fourier expansion. The key identity is stated in Equation~\eqref{eq:saw}.

\begin{equation}\label{eq:saw}
  S(x) = \frac{1}{2} - \frac{1}{\pi} \sum_{n=1}^{\infty} \frac{\sin(2\pi n x)}{n}
\end{equation}

\section{Convergence}
The partial sums converge for every $x \notin \mathbb{Z}$, with rate governed by \alpah.

\subsection{Partial sums}
% TODO: add figure of partial sums
Near each integer the partial sums overshoot the jump by a fixed proportion.

\section{The Gibbs phenomenon}\label{sec:gibbs}
% TODO: cite the original note
Near each jump the overshoot does not vanish as more terms are added.


\end{document}
